\documentclass{situationist-preprint}

\title{Living Without Illusions}
\msid{EXI-18-a}
\msfamily{EXI-18}
\msdiscipline{Meaning, Death and the Absurd}
\mstype{Principal manuscript}
\mskeywords{illusion; positive illusions; truth; evidence; self-deception; hope; contact with reality}

\begin{document}
\maketitle

\begin{abstract}
Can one live without illusions? Nietzsche held that false judgments may be indispensable to life, and a body of psychological research has suggested that mildly inflated self-assessments, an exaggerated sense of control and unrealistic optimism are features of normal functioning rather than of disorder. We argue that one can live without illusions only partly. Some illusions are psychologically necessary, in the sense that no one can function as an agent without some of them and their wholesale removal is not within anyone's power. But their necessity does not make them true, and it does not make them epistemically justified. Necessity is a fact about the believer; truth and justification depend on the world and on the evidence. Since belief is not under direct voluntary control, the necessity of an illusion gives no license to adopt it; and since contact with reality is part of well-being, it gives no reason to cultivate more of it than one must. Where the function of an illusion can be served by hope or by practical commitment, which do not misrepresent, those attitudes are to be preferred.
\end{abstract}

\section{The problem}

Nietzsche wrote that the falseness of a judgment is not yet an objection to it, and that the falsest judgments may be the most indispensable to life \citep[\S4]{Nietzsche1886}. Ibsen gave the thought dramatic form in the \enquote{life-lie} that sustains the characters of \emph{The Wild Duck}. A psychological literature has seemed to confirm it. \citet{Alloy1979} found that depressed subjects judged their control over outcomes more accurately than nondepressed subjects, who overestimated it. \citet{Taylor1988} argued that unrealistically positive self-evaluations, exaggerated perceptions of control and unrealistic optimism are characteristic of normal human thought and promote the capacity to care for others, to be happy and to do creative and productive work. \citet{McKay2009} asked which false beliefs could be products of well-functioning cognition, and concluded that positive illusions are the clearest candidates for misbelief that is adaptive as such.

The question is whether a person who values truth can live without illusions, and, if she cannot, what follows. We define an illusion as a belief that is false, or held with confidence greater than the evidence supports, in a direction favorable to the believer's interests or self-image.

\section{The argument}

\begin{definition}[Psychological necessity]\label{d:nec}
An illusion, or a class of illusions, is \emph{psychologically necessary} for an agent if she could not function as an agent, in the sense of forming and pursuing plans with normal persistence and without disabling distress, unless she had it, and if its removal is not within her power.
\end{definition}

\begin{proposition}[Partial necessity]\label{p:partial}
Some illusions are psychologically necessary for normal human agents; not all illusions are.
\end{proposition}

\begin{proof}
The first conjunct rests on the evidence just cited, read conservatively: the tendencies \citet{Taylor1988} describe are pervasive and are associated with persistence and well-being, and their thorough removal is not something anyone knows how to achieve. It is also supported by a structural point: since belief is not under direct voluntary control (EPI-11), an agent cannot remove an illusion by deciding to, and her indirect control through attention is limited by what she can bear to attend to. The second conjunct follows from the variation among people and among beliefs: many illusions are corrected by evidence without any loss of function, and the ordinary progress of a life consists partly in such corrections.
\end{proof}

\begin{theorem}\label{t:nottrue}
That an illusion is psychologically necessary does not make it true, and does not make it epistemically justified.
\end{theorem}

\begin{proof}
Whether a belief is true depends on whether things are as it represents them; whether it is justified depends on the evidence bearing on it. Psychological necessity, by Definition~\ref{d:nec}, is a fact about the agent's capacity to function. It is not a fact about the matter the belief concerns, and it is not evidence bearing on it: that I could not function without believing that my prospects are better than average is no indication that they are. So necessity confers neither truth nor justification.
\end{proof}

\begin{corollary}[No license]\label{c:license}
The necessity of an illusion does not license adopting it, and gives no reason to cultivate illusions beyond those one cannot avoid.
\end{corollary}

\begin{proof}
Adoption at will is not possible (EPI-11), and cultivation through selective attention would produce belief on insufficient evidence. Where the stakes of error are high and evidence is available, that is a culpable failure of inquiry (EPI-23). Where those conditions do not hold, what James's right licenses is a commitment to act on a hypothesis, not a credence that exceeds the evidence (EPI-23). Moreover, an illusion has a cost even when it helps: well-being includes contact with reality (ETH-05), and each illusion is a respect in which one's life is out of contact with how things are. So necessity explains why some illusions are unavoidable; it does not recommend more of them.
\end{proof}

\begin{remark}[Substitutes]\label{r:subst}
Several of the functions attributed to positive illusions can be served by attitudes that do not misrepresent. Hope sustains effort toward an outcome while being proportioned to its probability (EXI-15); a practical commitment to act as if a project will succeed sustains persistence without a belief that it will. To the extent that these attitudes serve the functions, the range of necessary illusions shrinks, and so does the cost of living with them.
\end{remark}

\section{Objections and replies}

\begin{objection}[Pragmatic justification]\label{o:prag}
If a belief is necessary for functioning, the agent has a decisive practical reason to have it, and a reason to have a belief is a justification of it.
\end{objection}

\begin{reply}
A practical reason to be in a state is not a reason that bears on whether the state's content is true. It may justify actions that make the state more likely, within the limits EPI-23 sets, but it does not make the belief justified as a belief. Theorem~\ref{t:nottrue} concerns epistemic justification only, and leaves room for the practical point.
\end{reply}

\begin{objection}[Sadder but wiser]\label{o:depr}
If accuracy goes with depression, the ideal of living without illusions recommends depression.
\end{objection}

\begin{reply}
It does not: the ideal is a value among others, and Corollary~\ref{c:license} counsels only against cultivating illusions beyond necessity, not against functioning. The empirical claim is also less secure than it once seemed; later work has qualified the finding that depressed judgment is generally more accurate. Our conclusion does not depend on it.
\end{reply}

\begin{objection}[Meaning needs illusion]\label{o:meaning}
Meaning depends on taking one's projects to matter more than any detached view can confirm. That is an illusion, and a necessary one.
\end{objection}

\begin{reply}
EXI-04 argues that the seriousness with which we hold our ends outruns any justification the detached standpoint must accept. That is not an illusion, since it involves no false belief: the serious agent need not believe that her ends are certified from the standpoint of the universe, only that they are worthwhile, which may be true (EXI-17). The fitting response to the discrepancy is irony, which sees it clearly.
\end{reply}

\section{Conclusion}

One can live without illusions only partly. Some illusions are psychologically necessary, in that normal agency depends on them and they cannot be removed at will. Their necessity is a fact about us, not about the world, and it makes them neither true nor justified. It explains the illusions we cannot shed; it does not license adopting or cultivating others, and where hope or practical commitment can do their work, those attitudes, which do not misrepresent, are better.

\begin{thebibliography}{9}\small
\bibitem[Alloy and Abramson(1979)]{Alloy1979} Alloy, L. B., and Abramson, L. Y. (1979). Judgment of contingency in depressed and nondepressed students: Sadder but wiser? \emph{Journal of Experimental Psychology: General}, 108(4), 441--485.
\bibitem[McKay and Dennett(2009)]{McKay2009} McKay, R. T., and Dennett, D. C. (2009). The evolution of misbelief. \emph{Behavioral and Brain Sciences}, 32(6), 493--510.
\bibitem[Nietzsche(1886)]{Nietzsche1886} Nietzsche, F. (1886). \emph{Jenseits von Gut und B\"ose}. Cited by section.
\bibitem[Taylor and Brown(1988)]{Taylor1988} Taylor, S. E., and Brown, J. D. (1988). Illusion and well-being: A social psychological perspective on mental health. \emph{Psychological Bulletin}, 103(2), 193--210.
\end{thebibliography}

\end{document}
